\documentclass[11pt,a4paper]{article}
\usepackage[T1]{fontenc}
\usepackage{lmodern,amsmath,amssymb,amsthm,mathtools,microtype,booktabs}
\usepackage[margin=27mm]{geometry}
\usepackage[colorlinks=true,linkcolor=blue,urlcolor=blue,citecolor=blue]{hyperref}
\setlength{\emergencystretch}{2em}
\newtheorem{theorem}{Theorem}[section]
\newtheorem{proposition}[theorem]{Proposition}
\newtheorem{lemma}[theorem]{Lemma}
\newcommand{\Q}{\mathbb Q}\newcommand{\Z}{\mathbb Z}
\DeclareMathOperator{\Tr}{Tr}
\title{Seven-contact detectors and a genus-twelve actual weight curve}
\author{}\date{17 November 2026}
\begin{document}\maketitle
\begin{abstract}
The ordinary-$abc$ derivative construction can use seven contact shifts
when its scale is changed from $M^4$ to $M^8$. All extraction, preparation
and finite-character guards transfer. The resulting complete eleven-variable
family has an independently checked degree bound $357309057$; this is an
arithmetic variant, not a degree record. Its sound nine-variable source has
physical lifting degree seven in the squared slope. In the enlarged auxiliary
point domain, the exact tuple $X=1$, $A=P$, $H=U=P/2$ has an actual weight
cover of genus twelve. The associated squared-slope and slope covers have
genera $90$ and $181$, respectively. These statements do not prove a guarded
rational point on either cover, nor the global realization needed for nine
variables. The general twenty-one-variable target is also unchanged.
\end{abstract}

\section{Source and arithmetic premises}
We retain ordinary rational $abc$ and exactly the elliptic group,
formal-comparison, trace-height, simultaneous-prime and intrinsic-pole
inputs of the derivative source \cite{Derivative}, with the auxiliary
point extension of \cite{FullAux}. No new assertion about those imported
foundations is made here. Put
\[
 K=\Q(\sqrt2),\quad E:y^2=x^3-(11-6\sqrt2)x,
 \quad P=\left(\frac{57-28\sqrt2}{8},\frac{-660+417\sqrt2}{32}\right).
\]
The retained group statement is $2E(K)=\Z P$. The root role $A=[n]P$
uses its original two-coordinate square-abscissa chart. The roles $H,U$
may use the guarded full affine curve, each with two rational abscissa
components and its finite ordinate algebra. The original finite point
algebra has rank $4^3=64$.

For $f(Q)=1/x(Q)$ and $\zeta(Q)=-x(Q)/y(Q)$ set
\[
 r=\frac{\zeta(U)}{\zeta(H)},\qquad
 r'=\overline{\frac{\zeta(U-A)}{\zeta(H-P)}},
\]
\[
 (C_1,\ldots,C_6)=
 \bigl(f(H),f(U),\overline{f(H-P)},\overline{f(U-A)},r-r',r-X\bigr).
\]
Every original point, ordinate, difference and quotient guard is retained;
individual tests $C_i$ may vanish. For a fixed odd integer $B\ge3$ put
\begin{equation}\label{eq:height}
 \beta=X^2+B,\quad T=\Tr x(A),\quad M=\beta^4+T^4,
 \quad a=M^8,\quad t=a^2-z^2=M^{16}-z^2.
\end{equation}
Here $z\in\Q$ is one free slope. Our seven shifts are
\[
 \mathcal S=\{-7,-6,-4,1,3,5,8\}.
\]
They are nonzero, distinct, sum to zero, and have no opposite pair.
Define
\begin{equation}\label{eq:weight}
 F(a,z)=\prod_{c\in\mathcal S}((a+c)^2-z^2),\qquad
 u=-14F/F_a,\qquad w=\frac{u}{(1+u)^2},
\end{equation}
where the derivative is taken before specializing $a$. The balanced
six-test detector is unchanged:
\begin{equation}\label{eq:detector}
 D_i=\frac{C_i^2}{1+uC_i^2},\quad
 W=\prod_{i=1}^6(1+D_i^2/u),\quad
 V=\Tr(W^{-1}),\quad b=\frac{4w}{V^2}.
\end{equation}
Retain $M,t,F,F_a,u,1+u\ne0$, all six balanced denominator norms,
the two-slot norm of $W$, the trace guard $V\ne0$, and the inherited
guards. The formerly required $2a+c+d\ne0$ follows here from $a>10$.
The existing fixed exceptional set $S_0$, containing every prime at most
$76176$, already contains the prime divisors of $14$, of every shift
difference, and of every nonzero pair sum. Thus no exceptional set changes.

The complete source retains two rational pole parameters $v,s$ and the
finite quadratic root
\begin{equation}\label{eq:pole}
 Y^2=\eta(\eta+v^2)^2s^4-\eta(\eta-v^2)^2(1-\eta)^2,
 \qquad \eta=7b.
\end{equation}
Only this final ordinate may vanish. The first six quadratic-stage
radicands in the point algebra remain units.

\section{Soundness with seven contacts}
For nonzero rational $q$ let $r^+_{\mathrm{odd},S_0}(q)$ be the product
of the primes outside $S_0$ having positive odd valuation on $q$.
The derivative-height theorem of \cite[Theorem 2.1]{Derivative} applies
to $a=M^m$ for every fixed $m$, provided the shifts have no opposite
pair. With seven shifts, $m=8$ and $\epsilon=1/100$ it gives
\begin{equation}\label{eq:radical}
 \frac{97}{200}h(t)
 \le \log r^+_{\mathrm{odd},S_0}(u)+\frac{201}{100}h(M)+O(1).
\end{equation}
This is precisely the existing one-variable $abc$ consequence; no
moving-form or multivariable version of $abc$ is being assumed.

The binary-power and signed mixed-power bounds \cite{Signed} give, for
$E_h=h(1:\beta:T)$ and $\kappa=1001/1000$,
\begin{equation}\label{eq:contractions}
 \alpha E_h\le h(M)+O(1),\qquad h(M)\le Qh(t)+O(1),
 \quad \alpha=\frac{997}{1001},\quad Q=\frac{501}{2995}.
\end{equation}
Indeed the signed mixed-power denominator is
$16-10\kappa=599/100>0$, and its numerator is
$2\kappa-1=501/500$. The precise absorption margin is
\begin{equation}\label{eq:margin}
 \frac{97}{200}-\frac{501}{2995}
       \left(\frac{201}{100}+\frac{1001}{1994}\right)
           =\frac{38695561}{597203000}>0.
\end{equation}

\begin{theorem}\label{thm:sound}
Under the retained premises, every guarded rational solution of
\eqref{eq:height}--\eqref{eq:pole} has $X=n\in\Z$ for every
sufficiently large fixed odd $B$.
\end{theorem}
\begin{proof}
The balanced valuation identity and the intrinsic pole implication are
unchanged. At a positive odd valuation $\nu=v_p(u)$ outside $S_0$,
write $q\le0$ for the largest valuation of $W$ in the two slots. The
absence of negative odd valuations of $b$ gives
$1\le v_p(b)\le\nu+2q$. Thus in a slot attaining $q$, each of the
six original tests has positive valuation. The formal-comparison theorem,
including its full-auxiliary extension, gives $p$-integrality of $X$ and
$X\equiv n\pmod p$. Consequently, if $X\ne n$,
\[
 \log r^+_{\mathrm{odd},S_0}(u)
       \le\tfrac12 E_h+\log(1+|n|).
\]
Combine this with \eqref{eq:radical}--\eqref{eq:contractions}.
The positive coefficient \eqref{eq:margin} yields
$h(t)\ll\log(1+|n|)+1$. The retained trace-height bound
$n^2\le A_0 h(T)+J_0$ and
$h(T)\le E_h\le Qh(t)/\alpha+O(1)$ then bound $n$ and $E_h$
independently of $B$. Since $E_h\ge h(\beta)\ge\log B$, this
is impossible for sufficiently large $B$. The root index remains
integral; allowing half-integral auxiliary indices does not alter the
formal comparison at any consumed odd prime.
\end{proof}

\section{Completeness and all guards}
\begin{theorem}\label{thm:complete}
The source has a guarded rational witness for every nonzero integer $X$.
For all sufficiently large fixed odd $B$, it therefore gives a complete
eleven-variable polynomial family, including the input.
\end{theorem}
\begin{proof}
Fix $X=n\ne0$, $A=[n]P$ and use the original subfamily $U=[n]H$
with $H\in2E(K)$. The enlarged auxiliary source contains this entire
subfamily. All original point and difference guards can therefore be
prepared and avoided exactly as in \cite{Derivative,FullAux}.

Since the shifts sum to zero,
\[
 F\sim-z^{14},\qquad F_a\sim14az^{12},\qquad
 u\sim z^2/a,\qquad w\sim a/z^2.
\]
The number $a=M^8$ is a rational square. Thus the fixed-place
preparation by making $z$ large again makes $u,w$ local squares
with the required even valuations, and sends $W\to1$, $V\to2$.
In particular it prepares the two-adic class $\eta\in7\Q_2^{\times2}$,
$v_2(\eta)\ge4$, and the place seven. At every prepared real
character choose $z>a+10$. Then $F<0$, $F_a>0$, $u,w>0$;
each balanced denominator is positive in both slots, and
$0<V\le2$. Taking $z$ larger makes $0<\eta<1$ at the selected
character. This also ensures every new denominator and trace norm
is nonzero on the entire real point-character locus used in completeness.

Choose $D>0$ with $D^2(a+c)\in\Z$ for $c\in\mathcal S$. The
same simultaneous-prime theorem now has fourteen distinct positions:
\[
 p_c^\pm=N\pm D^2(a+c)L^2,\qquad z=N/(D^2L^2).
\]
It retains the prescribed divisibility of $L$, order compatibility,
finite avoidance and real lower bound. Its coefficient set is fixed
once $n,A,B$ are fixed; no coefficient uniformity is required.
One has
\[
 F=-\frac{\prod_{c\in\mathcal S}p_c^-p_c^+}{D^{28}L^{28}}.
\]
At a selected prime just one linear factor has valuation one and
$F_a$ is a unit, so $v_p(u)=1$. The retained compatible index
congruences make the first four tests vanish to order at least two
and the comparisons to order at least four in a successful slot.
Therefore the original detector proof gives $V$ a unit and $v_p(b)=1$.

At unselected primes outside the temporary preparation, an integral
$z$ gives $v_p(u)\le0$. A pole of $z$ gives
$v_p(F)=14v_p(z)$ and $v_p(F_a)=12v_p(z)$, so $v_p(u)$ is even.
At the fixed numerator and denominator primes of $a$, sufficiently
strong leading-term dominance gives
$v_p(u)=2v_p(z)-v_p(a)$, also even because $v_p(a)$ is divisible by
eight. The same conclusion holds at primes dividing fourteen after
making the preparation sufficiently deep. At every negative odd
valuation of $u$, the unchanged balanced identity gives $W\equiv1$
in both slots, so $V\equiv2$ and $v_p(b)=-v_p(u)>0$.
Hence $\eta$ has square reduced denominator and satisfies all hypotheses
of the retained intrinsic-pole existence theorem.

That theorem supplies $v,s,Y$. Its Zariski-dense rational witness pairs
avoid the finite product of all nonzero character norms of the pole
and denominator guards, as before. No condition $\eta<1$ is required
at other characters; the polynomial in $v,s$ is nonzero at each character
because $\eta$ is a unit. Only the last root may be zero in soundness.
All point radicands and every preceding denominator retain their original
nonzero guards. This proves the full whole-character statement.
\end{proof}

The coordinate count is
$1+6+1+2+1=11$: input, point coordinates, slope, pole parameters,
and the single finite-algebra encoder. The integer circuit uses
$u=-7\mathcal G/\mathcal H$, where
$\mathcal G(a,Z)=\prod_c((a+c)^2-Z)$ and
$\mathcal H=\mathcal G_a/2$. The ordinary terminal-root construction
and its guarded subalgebra norms \cite{Terminal} give a single integer
polynomial with content one. The stored generator and independent
coefficient-DAG traversal certify 1940 nodes, 32 native guards, rank 128,
and the uniform total-degree upper bound
\[
                         357309057<2^{29}.
\]
The monomial bound is below $2^{270}$. For the stored fixture $B=3$,
the coefficient $\ell^1$ norm is at most $2^{727796736}$; this is a
bound for that fixture, not a sufficient defining value of $B$.
The generator computes a coefficient bound for every specified $B$.
The previous degree bound $232094337$ remains smaller.

\section{The nine-variable source and its unresolved requirement}
Replace \eqref{eq:pole} and its two free coordinates by one finite root
\begin{equation}\label{eq:actual}
                          Y^2=V^2+w(2-V).
\end{equation}
The local argument of \cite{Criterion} applies without alteration:
a negative odd valuation of $b$ outside $S_0$ would give
$v_p(V^2+w(2-V))=v_p(w)$ odd. Thus Theorem~\ref{thm:sound}
applies to this source as well. The two finite point and terminal-root
ranks remain $64$ and $128$; the coordinate count is nine.
All new rational denominators are already guarded by
\eqref{eq:weight}--\eqref{eq:detector}; no inverse of $2-V$ or of the
terminal ordinate is introduced.

Completeness requires actual prepared point tuples and a rational
slope giving a square in \eqref{eq:actual}, for every nonzero integer
input. No free trace or free weight may replace $V$ or $u$.
The fixed-place, selected-contact, negative odd valuation, and real
local-square statements of \cite{Criterion} transfer by the preceding
preparation proof. They do not establish a rational square.
The source is invariant under $(X,A,U)\mapsto(-X,-A,-U)$, because
only signs of the last two tests change. Hence realization for all
sufficiently large positive integers would suffice after multiplying
by finitely many integer input factors, without adding variables.
Neither this cofinite realization nor full nine-variable completeness
is proved here.

\section{An actual genus-twelve weight curve}
Take the algebraic fixture
\[
 X=1,\quad B=3,\quad A=P,\quad
 H=U=R_0:=(1+\sqrt2,2\sqrt2).
\]
Direct duplication gives $2R_0=P$. Thus $H-P=-H$ and $U-A=-H$.
The full auxiliary chart is essential here: $R_0$ is not being claimed
to belong to the old integral-index subfamily. Its affine abscissa,
ordinate, and required differences are nonzero; $x(R_0)\ne x(P)$.
Put $c=1/x(R_0)=\sqrt2-1$. The actual tests are
\[
                       (c,c,\bar c,\bar c,0,0).
\]
Define the following integer polynomials:
\begin{align*}
 D(u)&=u^2(u^2+6u+1)^2,\\
 R(u)&=D(u)+2u^3+12u^2+34u+1,\\
 \Pi(u)&=4D(u)^4(1+u)^2+2uR(u)^2\bigl(R(u)^2-D(u)^2\bigr).
\end{align*}
Here $D$ is a polynomial name, unrelated to the integer used in prime
preparation. For $\alpha=c^2=3-2\sqrt2$,
\[
 J(u)=\frac{\bigl(u(1+u\alpha)^2+\alpha^2\bigr)
              \bigl(u(1+u\bar\alpha)^2+\bar\alpha^2\bigr)}
             {u^2(1+6u+u^2)^2}=\frac{R(u)}{D(u)}.
\]
Consequently the detector and actual reciprocal trace are
\[
                  W=J(u)^2,\qquad V=\frac{2}{J(u)^2}
                                      =\frac{2D(u)^2}{R(u)^2}.
\]
This calculation uses the actual six tests, not independent constants.

\begin{proposition}\label{prop:small}
The actual weight cover is
\[
 Y_0^2=\Pi(u),\qquad Y_0=(1+u)R(u)^2Y.
\]
It is smooth after normalization and has genus twelve. At rational
weights, all selected-character balanced, product and trace denominator
guards hold whenever $u\ne0,-1$.
\end{proposition}
\begin{proof}
Substitute the displayed expression for $V$ into \eqref{eq:actual}
and clear the square denominator. The polynomial $\Pi$ has degree
26, leading coefficient four, and a simple zero at zero.
The certificate contains $e,f\in\mathbb F_{53}[u]$ satisfying
$e\Pi+f\Pi'=1$, with reduction preserving degree. Hence $\Pi$
is squarefree over $\Q$ and its normalization has genus $(26-2)/2=12$.

The quadratic $u^2+6u+1$ has nonsquare discriminant $32$, so it has
no rational zero. The monic polynomial $R$ has constant term one;
its only possible rational roots are $1,-1$, and
$R(1)=113$, $R(-1)=-7$. Thus $D,R$ are nonzero for rational $u\ne0$.
The norm of each balanced denominator is $1+6u+u^2$, and the product
and trace are nonzero by the formulas above. The guard $1+u\ne0$
is separate. These statements concern the selected actual character;
other characters and full preparation remain separate requirements.
\end{proof}

The branch point at $u=0$ and the two rational points at infinity lie
outside the guarded affine weight domain. No admissible rational point
is asserted by Proposition~\ref{prop:small}.

\section{The exact degree-seven physical lift}
For the fixed tuple above, the physical scale is
\[
                a=\left(\frac{10621537}{256}\right)^8.
\]
In this section only, let $T_0=z^2$. Form
\[
 F(a,T_0)=\prod_{c\in\mathcal S}((a+c)^2-T_0),\qquad
                   \frac{A(T_0)}{B(T_0)}=-14F/F_a,
\]
with $F_a$ differentiated before specializing $a$. Clear the common
rational denominators of this pair and divide its common integer
content. This defines the stored integral pair $A,B$ uniquely with
the indicated signs; its degrees are seven and six, respectively.
Set
\[
                P_0(T_0)=B(T_0)^{26}\Pi(A(T_0)/B(T_0)).
\]
It is kept as an exact homogeneous composition, not as an unnecessarily
large expanded list of coefficients.

\begin{proposition}\label{prop:physical}
The polynomial $P_0$ is squarefree of degree 182 and $P_0(0)\ne0$.
The normalized curves
\[
 Y_1^2=P_0(T_0),\qquad Y_2^2=T_0P_0(T_0),\qquad
                         Y_1^2=P_0(z^2)
\]
have genera $90,91,181$, respectively. A finite rational point
$(u,Y_0)$ of Proposition~\ref{prop:small} lifts to the actual source
exactly when it admits a rational $T_0$ with
\[
 uB(T_0)-A(T_0)=0,\quad B(T_0)\ne0,\quad
 T_0=z^2\ (z\in\Q),
\]
and all original physical guards. For full prepared realization, every
retained finite-place, contact, real, and whole-character requirement
must also be checked.
\end{proposition}
\begin{proof}
Modulo 53 the stored pair remains coprime and has degrees seven and
six. The homogeneous composition has degree 182 and nonzero constant
term. Its stored Bezout identity with its derivative certifies
squarefreeness over $\Q$. Multiplying by $T_0$ or substituting $z^2$
introduces no repeated root because the constant term is nonzero.
The hyperelliptic genus formulas give the three values.

The degree-seven equation is exactly the physical-weight relation,
and the separate square test is exactly rational recoverability of the
slope. The ordinate transformations are
$Y=Y_0/((1+u)R(u)^2)$ and $Y_1=B(T_0)^{13}Y_0$.
All denominators displayed here are guarded. Retain also
$a^2-T_0\ne0$ and every original point/chart guard. No implication
from quotient rationality to the missing lift or preparation is used.
\end{proof}

\section{Reproduction and scope}
The directory
\nolinkurl{computation/few_variables/seven_contact11_v74}
contains the complete circuit, its pinned dependencies and independent
bound traversal. Independent chord-and-tangent evaluations at two
finite fields check 384 sign characters, including half-integral
auxiliary indices. Deliberately corrupted degree and guard specifications
are rejected. These checks establish their stated finite assertions,
not the imported infinite arithmetic theorems.

The directory
\nolinkurl{computation/pole_parameters/halfpoint_trace_v74}
contains integer coefficient lists for $D,R,\Pi,A,B$, the two
finite-field Bezout certificates, and a second reconstruction that
imports no producer routine. Neither directory claims rational-point
enumeration, a prepared rational witness for the fixed tuple, or
nine-variable completeness. The old fixture $B=3$ is not certified to
exceed any sufficient defining threshold.

\begin{thebibliography}{9}
\bibitem{Derivative}\emph{An eleven-variable construction under rational $abc$},
\nolinkurl{papers/conditional_few_variables/derivative_detector11_20261117_v63.tex}.
\bibitem{FullAux}\emph{Full auxiliary elliptic points in the eleven-variable source},
\nolinkurl{papers/conditional_few_variables/full_auxiliary11_20261117_v71.tex}.
\bibitem{Signed}\emph{Signed mixed-power heights and single-cover criteria},
\nolinkurl{papers/elliptic_heights/signed_mixed_power_height_20261117_v58.tex}.
\bibitem{Criterion}\emph{One actual-trace square in a derivative detector},
\nolinkurl{papers/pole_parameters/single_trace_derivative_criterion_20261117_v64.tex}.
\bibitem{Terminal}\emph{A terminal quadratic root and complete polynomial specifications},
\nolinkurl{papers/finite_algebra/terminal_quadratic_root_20261115_v40.tex}.
\end{thebibliography}
\end{document}
